% v10 structural prefix-closure review, 9 October 2026.
% Candidate consolidated TODO after Definition 6.29 (tree of partial types).
% Review only; the author-level four mathematical repairs stay untouched.
%
\todo[inline]{\textbf{v10 validation (shortened partial-type realization).}
The one-filler construction now yields an actual
enumerated $L^+$-partial structure in the explicit
finite-unary/binary model: retain the first $d$ vertices,
insert a neutral filler at $d$, and relocate the type
vertex to $d+1$. Exact copying of the whole $L$-reduct
and the prescribed $E$ relation gives all three
partial-structure axioms, including the condition
that every nonempty directed $L$-pair has the required
$E$-incidence. The new ordinary vertex has
$\operatorname{fl}=d$, and its complete extracted
$L^+$-partial type agrees with
$\operatorname{type}_{A^+}^{d}(v)$, including E,
unary, diagonal, directed binary and negative facts.
This is formalised by
\texttt{PrefixReplicaPartial.lean}, with
\texttt{PrefixReplicaE.lean} and
\texttt{PrefixReplicaL.lean}.
To conclude closure of the \emph{admissible}
$K^\mathcal F_{pt}$ family under prefixes,
it remains to check that the replica's $L$-reduct
avoids the forbidden family. Use irreducibility:
a nontrivial forbidden structure cannot contain
the isolated filler, while the other vertices
project back into $A$ as an induced ordered copy;
a forbidden singleton at the filler is excluded
by the chosen allowed neutral singleton type.
Keep the final $K^\mathcal F_{pt}$ and Lemma~6.49
validation marks pending until this last age proof
and the actual forbidden-free family adapter.}
%
% New endpoint checks (green only on exact CI head):
%   prefixReplicaPartialStructure_freeLevel
%   prefixReplicaPartialStructure_type_eq
